\chapter{Conclusions}
\label{ch:conclusions}

\section{Summary}

This thesis set out to close a specific gap: the Standard Model Extension
(SME) and the Dobrescu--Mocioiu (DM) classification of exotic
spin-dependent potentials describe the same physics in two languages, one
relativistic and field-theoretic, the other non-relativistic and directly
tied to laboratory observables, with no complete, systematic dictionary
between them and no framework for using such a dictionary to compare
matter- and antimatter-sector constraints as a test of CPT symmetry.
Chapters~2 and~3 built the theoretical bridge, deriving the
non-relativistic Hamiltonians generated by $b_\mu$, $H_{\mu\nu}$, and
$d_{\mu\nu}$ via the Foldy--Wouthuysen transformation and matching each
onto the DM potential catalogue $V_1$--$V_{16}$. Chapter~4 applied that
bridge to real data, compiling 283 experimental constraint datasets (247 analysed) into a
single standardised database and computing the CPT asymmetry parameter
$A_\alpha$ for every matter--antimatter pair the database could support,
before turning the same database around to identify where it is currently
silent and proposing concrete next steps. The result is not a single
number or a single discovery, but a working pipeline: a mapping table
connecting relativistic coefficients to measurable potentials, a
standardised constraint database built to be extended, and a quantitative
account of exactly how far current data can and cannot go toward testing
CPT symmetry in the spin-dependent sector.

\section{Revisiting the Objectives}

All five specific objectives set out in Chapter~1 were achieved in full.
The SME--DM mapping was derived completely for $b_\mu$, $H_{\mu\nu}$ and
$d_{\mu\nu}$, reproducing the expected operator structure and CPT
behaviour at every order in $1/m$ kept in this thesis. All three
$d_{\mu\nu}$ components, $d_{i0}\to V_2$, $d_{ij}\to V_7,V_8$ and
$d_{00}\to V_8$, match Kostelecký and Lane (1999) exactly, including their
signs, once that reference's momentum-index convention is applied
consistently; an apparent sign discrepancy in the two momentum-dependent
components was traced to that convention rather than to the derivation. A standardised database
of 283 datasets was compiled across five experimental platforms, 247 of
which entered the analysis. The CPT
asymmetry parameter was defined and computed for all fifteen matchable
matter--antimatter pairs, with the important qualification, itself one of
this thesis's principal findings, that $A_\alpha$ computed from one-sided
experimental bounds cannot by itself distinguish genuine CPT violation
from an ordinary sensitivity gap. The antimatter-sector coverage gap was
quantified by both fermion sector and DM potential, and concrete
experimental and archival strategies were proposed for closing it.

\section{Principal Findings}

The master matching table (Table~\ref{tab:master}) is this thesis's
central technical deliverable. Its single most physically significant
result is the $b_i$ derivation: a genuinely CPT-odd coefficient produces an
exact sign flip between matter and antimatter, obtained independently from
the Lagrangian's general CPT transformation and from an explicit
hole-theory derivation of the antiparticle Hamiltonian, after a
charge-conjugation argument on the bilinear alone was checked and found
insufficient to produce the sign flip on its own (Section~\ref{sec:results-fw-bmu}).
The most conceptually important result,
however, is methodological rather than numerical: the asymmetry parameter
$A_\alpha=(g_f-g_{\bar f})/(g_f+g_{\bar f})$, as actually computed from
published experimental limits, cannot distinguish a genuine CPT-violating
signal from an ordinary difference in sensitivity between the matter- and
antimatter-sector experiments being compared. An exact signed CPT-odd
relation makes the formula diverge rather than saturate at $1$; what
actually drives real output to $|A_\alpha|\approx1$ is the ordinary
behaviour of a ratio of two independent, always-positive upper bounds of
very different tightness, a pattern that arises regardless of the true CPT
parity of the underlying physics. This was confirmed empirically: the
CPT-odd $g_Ag_A$ pairs and the CPT-even $g_sg_s$ pair alike showed
statistically significant asymmetries after correcting for the correlated
nature of interpolated constraint points, yet none can currently be read
as evidence of CPT violation rather than sensitivity-gap saturation. Of the
247 datasets analysed, only 19 involve an antimatter sector, and from
these only fifteen matter--antimatter pairs could be matched; seven of twelve
classified DM potentials and six of eleven fermion-sector pairs have zero
antimatter-sector coverage despite substantial matter-sector data. This is
the binding constraint identified by this thesis: not a lack of
matter-sector precision, but an almost complete absence of comparably
precise antimatter-sector measurements in exactly the channels where
matter-sector data is most abundant.

\section{Limitations}

Two limitations bear on how much weight the results above can be given.
Thirty-six of the
283 compiled datasets (roughly one in eight) are held out of the
analysis rather than classified, twelve of them source-flagged
\texttt{review\_only} envelope curves to which no single potential applies;
and a small number of entries report physically implausible interaction
ranges likely caused by a unit-detection error. Neither affects the fifteen
matched pairs analysed, but the unit error should be resolved before a
broader coverage claim is drawn from the full database. Finally,
this thesis derives the non-relativistic limit of only three of the eight
SME coefficients classified in Chapter~2; the remainder are discussed
qualitatively but not derived with the same rigour and are not represented
in the asymmetry analysis.

\section{Recommendations for Future Work}

The most direct theoretical extension is to apply the Foldy--Wouthuysen
machinery, which is general and applies unchanged, to the remaining SME
coefficients $a_\mu$, $c_{\mu\nu}$, $e_\mu$, $f_\mu$, and
$g_{\lambda\mu\nu}$. The kinetic-sector field redefinition already required
for $d_{\mu\nu}$ carries over directly to $c_{\mu\nu}$, which occupies the
same position in the modified kinetic operator, so that coefficient is the
natural next one to derive. On the database side, decomposing the twelve
$g_Ag_V$ \texttt{Combined} envelope curves into their constituent
per-potential bounds, and settling the one unassigned astrophysical bound
against its source publication, are concrete and bounded cleanup tasks, and
the populated but unmatched $e$-$\mu$/$e$-$\bar\mu$ sector is a small,
well-defined follow-up that could add a twelfth matched pair without new
experimental input. Experimentally, new
antimatter-sector bounds should be prioritised in the $e$-$N$, $n$-$N$,
$p$-$N$, $e$-$n$, $n$-$p$, and $n$-$n$ sectors and in the $V_{1a}$, $V_{9+10}$, $V_1$, and
$V_{12+13}$ potentials, and existing BASE, ALPHA, and co-magnetometer
publications should be searched for data that could be reinterpreted
within this framework before new measurements are commissioned. Throughout,
a new antimatter-sector bound is useful for a CPT test specifically when it
approaches the precision of its matter-sector counterpart, not merely when
it exists.

\section{Closing Remarks}

The gap this thesis set out to address, a missing dictionary between the
SME and the Dobrescu--Mocioiu potentials, and no systematic way to use that
dictionary to compare matter- and antimatter-sector constraints, is now
substantially, though not completely, closed. The mapping table, the
compiled 283-dataset database, and the gap analysis together give a
reproducible, extensible pipeline rather than a fixed set of numbers: the
database is built to absorb new datasets as they are published, and the
pairing and asymmetry-analysis methodology apply unchanged to any new
antimatter-sector measurement that lands in one of the gap channels
identified here. The most durable contribution of this thesis may not be
any single derived Hamiltonian or compiled bound, but the demonstration
that an asymmetry parameter built from one-sided experimental bounds is a
sensitivity-gap diagnostic before it is a CPT diagnostic, and that closing
the antimatter-sector coverage gaps mapped in this thesis is a
precondition for that parameter to mean what it was designed to mean.
